\subsection{Training objectives}
\label{sec:training}

\paragraph{Noise and Data Interpolation via Minibatch coupling.}
Given minibatches \( x^0 \sim \pi_0 \) and \( x^\mathrm{data} \sim \pi_\mathrm{data} \) where $\pi_0$ is the uniform noise and $\pi_{data}$ is the data distribution respectively. We pair samples using a permutation-invariant approximation of \( c_{\mathrm{FGW}}(x^0,x^\mathrm{data}) \) described in \Cref{sec:gem}, via histogram matching (\Cref{app:matching}), and solving for a transport map \( T \) induced by the minibatch coupling, using the \texttt{POT} solver \citep{Flamary2021}.
 Each source graph $x^0$ is paired with a data graph $x^\mathrm{data}=T(x^0)$.
% In all experiments involving our method, we initialize from a uniform noise distribution $\pi_0$.

Crucially, we utilize $c_{\mathrm{FGW}}$ solely to select optimal pairs based on global structural similarity (e.g., node and edge counts). %Once paired, we rely on the isotropy of the uniform source: since the noise graph lacks intrinsic node ordering, any node permutation is statistically equivalent. We therefore treat the sampled node ordering of $x^0$ as implicitly aligned with the paired $x^{\mathrm{data}}$ via the identity permutation. 
% This approach is compatible with our architecture because $V_\theta$ is permutation-invariant by construction (\Cref{app:perm-invariance}), ensuring the learned energy landscape is independent of the specific stochastic node labeling.
Denote their embeddings by $\hat{x}^0$ and $\hat{x}^\mathrm{data}$. Define the
displacement in the embedding space as $v=\hat{x}^\mathrm{data}-\hat{x}^0$.
We sample a \emph{discrete} interpolant $x^t:=(x^t_\mathcal{V},x^t_\mathcal{E})$, where $x_\mathcal{V}^t:= [x_i^t]_{i=1}^n$ and $x_\mathcal{E}^t := [x_{ij}^t]_{i<j}$ by independently sampling each node and edge from the endpoint graphs.
\begin{eqnarray}
x_i^t \sim \operatorname{Cat}\Big((1-t)\,\phi_\mathcal{V}\big(x_i^0\big) + t\,\phi_\mathcal{V}\big(x_i^{\mathrm{data}}\big)\Big)\\
x_{ij}^t \sim \operatorname{Cat}\Big((1-t)\,\phi_\mathcal{E}\big(x_{ij}^0\big) + t\,\phi_\mathcal{E}\big(x_{ij}^{\mathrm{data}}\big)\Big)
\end{eqnarray}
% \end{equation}
where $\operatorname{Cat}$ denotes categorical sampling. Let $\hat{x}^t$ be the continuous embedding of the discrete graph $x^t$ via $\phi_\mathcal{V},\phi_\mathcal{E}$ (\Cref{sec:gem}). While this embedding resides in a continuous ambient space, training is conducted solely on discrete graphs $x \in \mathcal{X}$, as sampling proposals are concerning discrete jumps. The continuous embedding thus mainly facilitates gradient computation of $\nabla_{\hat{x}_t} V_\theta(\hat{x}^t)$ to guide these discrete transitions.



\paragraph{Objectives.}
We use the optimality condition derived in \Cref{proposals} to construct a flow-like loss ($\mathcal{L}_{\text{Flow}}$). We set the JKO step-size parameter $\eta=1$. Consequently, the optimality condition $\nabla_{\hat{x}} \tilde{V_\theta}(\hat{x}) = - \frac{1}{2 \eta}\nabla_{\hat{x}} \tilde{c}(\hat{x},\hat{y})$ from \eqref{eq:optimality} implies that the negative energy gradient should match the displacement vector $v$ pointing towards the data. %While the pairing is determined by global signatures,
The training objective minimizes the cost ($c_{\text{loc}}$) along the interpolation path:
\begin{align}
\mathcal{L}_{\text{Flow}}(\theta)
&=
\mathbb{E}_{(x^0,x^\mathrm{data})\sim\tilde{\gamma},\,t\sim\Unif([0,1])}\left[
\|\nabla_{\hat{x}_t} \tilde{V}_\theta(\hat{x}^t)+v\|_2^2
\right],\\
\mathcal{L}_{\text{CL}}(\theta)
&=
\mathbb{E}_{x^{+}\sim \pi_\mathrm{data}}[\tilde{V}_\theta(\hat{x}^{+})]
-\mathbb{E}_{x^{-}\sim \text{sg}\,(\pi_\theta)}[\tilde{V}_\theta(\hat{x}^{-})].
\label{eq:cl}
\end{align} 
The contrastive loss $\mathcal{L}_{\text{CL}}$~\cite{hinton2002training} (see \Cref{app:energy-refinement} for the derivation) leverages the discrete proposals introduced in \Cref{proposals} and \Cref{proposals_lang} to efficiently sample from the model distribution.
 We minimize:
\begin{equation}
\min_\theta\;
\mathcal{L}_{\text{Flow}}(\theta)
+\lambda_{\text{CL}}\,
\mathcal{L}_{\text{CL}}(\theta).
\end{equation}


We perform a warm-up phase of $N_{\mathrm{warmup}}$ training iterations using only $\mathcal{L}_{\text{Flow}}$ before introducing $\mathcal{L}_{\text{CL}}$. This serves two purposes: first, consistent with the intuition from the continuous formulation~\citep{balcerak2025energy}, it yields higher-quality negatives for contrastive learning; second, it allows us to \emph{temporarily} freeze the network after warm-up, solely to calibrate the mixing-phase sampler hyperparameters $(\beta_{\mathrm{mh}}, \beta^L, \lambda_{\mathcal V}^L, \lambda_{\mathcal E}^L)$ required by $\mathcal{L}_{\text{CL}}$. Specifically, we tune these hyperparameters to minimize the energy of samples generated by the fixed $V_\theta$, which is equivalent to maximizing their relative likelihood.

To approximate samples from \( \pi_\theta \) required by \(\mathcal{L}_{\text{CL}}\), we run Markov chains initialized in equal proportions from uniform noise and data samples.
(See \Cref{app:exp-setup} for detailed hyperparameters.) We outline high-level algorithms:
\begin{algorithm}[H]
\caption{GEM Training}
\small
\begin{algorithmic}[1]
\STATE {\bfseries Input:} $N_{\text{CL}}, N_{\text{warmup}}, \lambda_{\text{CL}}$, step $i=0$
\WHILE{$V_\theta$ not converged}
    \STATE Sample $(\xzero,\xdata)\sim\pizero\times\pidata$ (via SOLVER), $t\sim\mathcal{U}(0,1)$
    \STATE $x^t\sim \operatorname{Cat}\big((1 - t)\hat{x}^0 + t\hat{x}^{\text{data}}\big)$
    \STATE $\mathcal{L}_{\text{Flow}}=\|\nabla\tilde{V}_\theta(\hat{x}^t)+(\hat{x}^{\text{data}}-\hat{x}^0)\|^2$
    \IF{ $i>N_{\text{warmup}}$}
        \STATE $x \sim \text{GEM Sampling}(V_\theta) \hfill \leftarrow\ \text{Algorithm~2}$
        \STATE $\mathcal{L}_{\text{CL}}=\mathbb{E}[\tilde{V_\theta}(\hat{x}^{\text{data}})]-\mathbb{E}[\tilde{V_\theta}(\hat{x})]$
    \ELSE
        \STATE $\mathcal{L}_{\text{CL}}=0$
    \ENDIF
    \STATE \textbf{optimizer\_step}($\mathcal{L}_{\text{Flow}}+\lambda_{\text{CL}}\mathcal{L}_{\text{CL}}$) 
    \STATE $i=i+1$
\ENDWHILE
\end{algorithmic}
\end{algorithm}

\begin{algorithm}[H]
\caption{GEM Sampling}
\small
\begin{algorithmic}[1]
\STATE {\bfseries Input:} $V_\theta$, $N_{\text{CL}}$, $V_\text{th}=\mathbb{E}_{\pidata}[\Vtheta]$
\STATE Init. $x\sim\pizero$ (or $\pidata$)
\WHILE{$\Vtheta(x)>V_\text{th}$ {\bfseries and} $x$ not stuck}
    \STATE $x\sim \qgreedy(x\to y)$
\ENDWHILE
\FOR{$k=1$ {\bfseries to} $N_{\text{CL}}$}
    \STATE $x\sim \text{MH-step}(\qmixing(x\to y))$
\ENDFOR
\STATE {\bfseries Return} $x$
\end{algorithmic}
\end{algorithm}
\vspace{-1em}